\section{Induction}

\subsection{Linear categories and functors}

We will define here linear category theory.
In the language of category theory, this is simply the theory of categories enriched over the monoidal category of vector spaces (over a given field).
But we will ignore this abstraction, and instead define everything from the ground-up.

\bnote

    I have grown weary of denoting fields using blackboard bold (e.g.\ $\bF$).
    Instead a typical field will be denoted by $k$.
    Blackboard bold seems to give more significance to the field than necessary in my opinion.

\enote

\bdefn

    A {\emphcolor $k$-linear category} (or just a category) $\C$ consists of:
    \benum
        \item a class of objects (we abuse notation and denote this class also by $\C$);
        \item for every two objects $M,N\in\C$, a $k$-vector space $\hom_\C(M,N)$, called a {\emphcolor hom-space};
        \item for every three objects $M,N,L\in\C$, a $k$-bilinear map
        $$ {\circ}\colon\hom(N,L)\times\hom(M,N)\to\hom(M,L) $$
        (Equivalently, this is a $k$-linear map from the tensor product of the left hom-spaces.)
        Such that for every $M\in\C$ there exists a ${\rm id}_M\in\hom(M,M)$ which is the identity under composition (left/right according to where composition is defined), and $\circ$ is associative
        when defined.
    \eenum

\edefn

When declaring $\alpha\in\hom_\C(M,N)$ we may instead write $\alpha\colon M\to N$.

For example, given a group $G$ and field $k$, we can form a category $\Rep(G)$ whose objects are $G$-representations over $k$, and morphisms are $G$-equivariant morphisms.
In particular, the hom-spaces are simply $\hom_G(M,N)$.

Similarly, for a $k$-algebra $A$, we can consider the category ${}_A\Mod$ of left $A$-modules, and $\Mod_A$ of right $A$-modules.

\bdefn

    Let $\C$ and $\D$ be $k$-linear categories.
    A {\emphcolor functor} from $\C$ to $\D$, denoted $F\colon\C\to\D$, associates for every object $M\in\C$ an object $F(M)\in\D$, and for every $M,N\in\C$ a $k$-linear map
    $$ F\colon\hom_\C(M,N) \to \hom_\D(FM,FN) $$
    such that $F{\rm id}_M={\rm id}_{FM}$ and $F(\alpha\circ\beta)=F(\alpha)\circ F(\beta)$.

\edefn

\bdefn

    Let $F,G\colon\C\to\D$ be functors.
    A {\emphcolor natural transformation} $\sigma\colon F\to G$ consists of associating to every $M\in\C$ a map $\sigma_M\colon FM\to GM$ such that for every $\alpha\colon M\to N$ in $\C$,
    we have
    $$ G(\alpha)\circ\sigma_M = \sigma_N\circ F(\alpha) $$
    (Write out the commutative square for clarity.)

\edefn

A natural transformation for which every component ($\sigma_M$) is an isomorphism (in the usual category-theoretic sense) is called a {\it natural isomorphism}.
Two functors are {\it naturally isomorphic} if there exists a natural isomorphism between them.

Note that we can easily define the composition of functors, the product of categories and functors, opposite categories, in the standard way.
In particular we can define the {\it hom-functor}:
$$ \hom\colon\C^\op\times\C\to{\rm Vec}_k $$
which maps pairs of objects $(M,N)$ to $\hom_\C(M,N)$ and pairs of morphisms $f\colon M\to N$ (in $\C$, so $N\to M$ in $\C^\op$) and $g\colon L\to K$ to
$$ \hom(f,g)\colon\hom(N,L)\to\hom(M,K),\qquad \hom(f,g)(\alpha) = g\circ\alpha\circ f $$
For a set object $M\in\C$, this restricts to give two functors:
$$ \hom(M,\bul)\colon\C\to{\rm Vec}_k,\qquad\hom(\bul,M)\colon\C\to{\rm Vec}_k $$
where $\hom(M,f)$ is post-composition with $f$ (i.e.\ $\hom(M,f)\alpha=f\alpha$) and $\hom(f,M)$ is pre-composition with $f$ ($\hom(f,M)\alpha=\alpha f$).

We now define adjoint pairs;

\bdefn

    Let $F\colon\C\to\D$ and $G\colon\D\to\C$ be two functors.
    They form an {\emphcolor adjoint pair} (or an {\emphcolor adjunction}) if the functors $\hom_\D(F\bul,\bul)$ and $\hom_\C(\bul,G\bul)\colon\C^\op\times\D\to{\rm Vec}_k$ are naturally isomorphic.
    In such a case, we write $F\dashv G$, and call $F$ the {\emphcolor left adjoint} and $G$ the {\emphcolor right adjoint}.

\edefn

The naturality condition is equivalent to saying that for every $M\in\C$ and $U\in\D$ there exists a $k$-linear isomorphism $\gamma_{MU}\colon\hom_\D(FM,U)\to\hom_\C(M,GU)$ such that
for every $\alpha\colon M_1\to M_2$ in $\C$ the following diagram commutes:

\centerline{\drawdiagram{
    $\hom_\D(FM_2,U)$&$\hom_\C(M_2,GU)$\cr
    $\hom_\D(FM_1,U)$&$\hom_\C(M_1,GU)$
}{
    \diagarrow{from={1,1}, to={1,2}, text=$\gamma_{M_2U}$, y distance=.25cm}
    \diagarrow{from={2,1}, to={2,2}, text=$\gamma_{M_1U}$, y distance=.25cm}
    \diagarrow{from={1,1}, to={2,1}, text=$\bul\circ F\alpha$, x distance=-.5cm}
    \diagarrow{from={1,2}, to={2,2}, text=$\bul\circ \alpha$, x distance=.5cm}
}}

as well as for every $\beta\colon U_1\to U_2$ in $\D$ an analogous diagram commute.

An equivalent definition of adjunctions are via unit-counit pairs.

\bdefn

    Let $F\colon\C\to\D$ and $G\colon\D\to\C$ be two functors.
    A {\emphcolor unit-counit pair} for $(F,G)$ (note the order) are two natural transformations $\eta\colon{\rm Id}_\C\to GF$ (the unit) and $\epsilon\colon FG\to{\rm Id}_\D$ such that the compositions
    $$ F\xvarrightarrow{F\eta}FGF\xvarrightarrow{\epsilon F}F,\qquad G\xvarrightarrow{\eta G}GFG\xvarrightarrow{G\epsilon}G $$
    are the respective identity natural transformations.

\edefn

In the above definition, we use the concept of {\it whiskering}.
For a natural transformation $\sigma\colon F\to G$ and another functor $H$, we define the natural transformation $H\sigma\colon HF\to HG$ (if such composition is well-defined) by $(H\sigma)_X=H(\sigma_X)$
and $\sigma H\colon FH\to GH$ by $(\sigma H)_X=\sigma_{H(X)}$.
So the unit-counit conditions say that
$$ {\rm id}_{FX} = \epsilon_{FX}\circ F(\eta_X),\qquad {\rm id}_{GY} = G(\epsilon_Y)\circ\eta_{GY} $$
for all $X\in\C$ and $Y\in\D$.

\bthrm

    $F\dashv G$ form an adjunction if and only if there exists a unit-counit pair $(\eta,\epsilon)$ for $(F,G)$.

\ethrm

\bproof

    Given an adjunction $F\dashv G$, let
    $$ \eta_X=\gamma_{X,FX}({\rm id}_{FX})\in\hom_\D(FX,FX),\quad\hbox{and}\quad\epsilon_Y=\gamma^{-1}_{GY,Y}({\rm id}_{GY})\in\hom_\C(GY,GY) . $$
    This forms the unit-counit pair.

    Conversely, given a unit-counit pair $(\eta,\epsilon)$, define $\gamma\colon\hom_\D(F\bul,\bul)\to\hom_C(\bul,G\bul)$ by $\gamma_{XY}(\alpha)=G(\alpha)\circ\eta_X$ for $f\in\hom_\D(FX,Y)$.
    Its inverse is $\gamma^{-1}_{YX}(\beta)=\epsilon_Y\circ F(\beta)$.
    \qed

\eproof

Let's construct an example of an adjunction.
Let $\iota\colon B\to A$ be a morphism of $k$-algebras.
Then every $A$-module can be viewed as a $B$-module; given an $A$-module $M$, give it a $B$-module structure by $b\cdot m=\iota(b)m$ for $b\in B,m\in M$ (this is just composition of $\iota$ with the $A$-module
structure).
Note that with such a view, every $A$-module morphism is a $B$-module morphism: $f(bm)=f(\iota(b)m)=\iota(b)f(m)=bf(m)$.
Thus we can define a functor
$$ \res^A_B\colon\Mod(A)\to\Mod(B) $$
which maps an $A$-module $M$ to its induced $B$-module structure, and an $A$-module morphism to itself.
$\res^A_B$ is called the {\it restriction functor} (or the {\it forgetful functor}).

Now, we ask ourselves if the forgetful functor has a left adjoint.
Consider the functor
$$ A\otimes_B\bul\colon\Mod(B)\to\Mod(A) $$
which maps a $B$-module $M$ to the $A$-module $A\otimes_BM$ (this is well-defined, $A$ is a right $B$-module by $a\cdot b=a\iota(b)$).
And it maps a morphism of $B$ modules $f\colon M_1\to M_2$ to $A\otimes f\colon A\otimes_BM_1\to A\otimes_BM_2$ which maps $a\otimes m$ to $a\otimes f(m)$.
Such a morphism exists and is unique.

We claim that $A\otimes_B\bul$ and $\res^A_B$ forms an adjunction.
Given $M\in\Mod(B)$ and $N\in\Mod(A)$ we define an isomorphism
$$ \gamma_{MN}\colon\hom_{\Mod(A)}(A\otimes_BM,N) \to \hom_{\Mod(B)}(M,\res^A_BN) $$
by mapping $\alpha\colon A\otimes_BM\to N$ to $\beta\colon M\to\res^A_BN$ defined by $\beta(m)=\alpha(1\otimes m)$.
Its inverse,
$$ \gamma_{MN}^{-1}\colon\hom_{\Mod(B)}(M,\res^A_BN) \to \hom_{\Mod(A)}(A\otimes_BM,N) $$
is given by mapping $\beta\colon M\to\res^A_BN$ to $\alpha\colon A\otimes_BM\to N$ by $\alpha(a\otimes m)=\iota(a)\beta(m)$.
Their inverse-ness is readily checked, as well as their naturality.

Thus we have
$$ A\otimes_B\bul\,\dashv\,\res^A_B $$

Interestinly, $\res^A_B$ has a right adjoint as well.
Namely, consider the functor $\Mod(B)\to\Mod(A)$ $M\mapsto\hom_B(A,M)$ ($A$ is a $B$-module via $\iota$).
$\hom_B(A,M)$ is a $A$-module by $(a\phi)(a')=\phi(a'a)$.
Given a morphism $\alpha\colon M_1\to M_2$ we define $\alpha_*\colon\hom_B(A,M_1)\to\hom_B(A,M_2)$ by post-composition.

Now given $M\in\Mod(B),N\in\Mod(A)$ we define
$$ \gamma_{MN}\colon \hom_{\Mod(B)}(\res^A_BN,M) \to \hom_{\Mod(A)}(N,\hom_B(A,M)) $$
by mapping $\alpha\colon\res^A_BN\to M$ to $\beta\colon N\to\hom_B(A,M)$ defined by $\beta(n)(a)=\alpha(an)$.
Its inverse maps $\beta\colon N\to\hom_B(A,M)$ to $\alpha(n)=\beta(n)(1)$.
All the conditions are readily checked.

Thus we have a chain of adjunctions:
$$ A\otimes_B\bul\,\dashv\,\res^A_B\,\dashv\,\hom_B(A,\bul) $$
when $\iota\colon B\to A$ is a $k$-algebra morphism.

\subsection{Basics of induction}

Fix groups $G,H$, a base field $k$, and a group morphism $\iota\colon H\to G$.

$\iota$ induces a morphism of $k$-algebras $\iota\colon k[H]\to k[G]$.
Since there is an isomorphism of categories $\Rep_k(G)\cong\Mod(k[G])$, we have a chain of adjunctions:

\kern1cm
\centerline{\drawdiagram{
    $\Rep(H)$&$\Rep(G)$
}{
    \diagarrow{from={1,1}, to={1,2}, text=$\ind^H_G$, y off=.5cm, y distance=.75cm}
    \diagarrow{from={1,1}, to={1,2}, text=$\Ind^H_G$, y off=-.5cm, y distance=-.25cm}
    \diagarrow{from={1,2}, to={1,1}, text=$\res^G_H$, y distance=.25cm}
}}
\kern.5cm

Where
$$ \ind^H_G\dashv\res^G_H\dashv\Ind^H_G $$
Specifically $\res^G_H=\res^{k[G]}_{k[H]}$, $\ind^H_G$ is $k[G]\otimes_{k[H]}\bul$, and $\Ind^H_G$ is $\hom_{k[H]}(k[G],\bul)$.

\subsubsection*{Understanding $\ind^H_G$}

We will assume for now that $\iota\colon H\to G$ is an inclusion.
We will examine first $\ind^H_G$.
Let $(g_q)_{q\in G/H}$ be a set of representatives of the cosets in $G/H$ (so $g_q\in q$).
Then we have an isomorphism of right $k[H]$-modules:
$$ k[G]\cong\bigoplus_{q\in G/H}k[H] $$
given by mapping $(d_q)_{q\in G/H}$ on the right to $\sum_{q\in G/H}\delta_{g_q}\cdot d_q$ on the left.
It is readily checked that this is an injective morphism of right $k[H]$-modules, and by considering dimensions we have that it is an isomorphism.
Thus, given a (left) $k[H]$-module $M$:
$$ k[G]\otimes_{k[H]}M \cong \parens{\bigoplus_{q\in G/H}k[H]}\otimes_{k[H]}M \cong \bigoplus_{q\in G/H}(k[H]\otimes_{k[H]}M) \cong \bigoplus_{q\in G/H}M $$
Sending $(m_q)_{q\in G/H}$ on the right to $\sum_{q\in G/H}\delta_{g_q}\otimes m_q$ on the left.

So we can think of $\ind^H_GM$ as $\bigoplus_{q\in G/H}g_qM$ (where $g_q$ ``acts'' on $M$ to give $M$), and we have $g(g_qm)=g_{q'}(hm)$ for $g\in G$ where $gg_q=g_{q'}h$ for $h\in H$.
In particular, we have
$$ \dim(\ind^H_G(V)) = [G:H]\cdot\dim V $$

So suppose we have a $G$-representation $V$, and an $H$-subrepresentation $W\subseteq V$.
Then as the inclusion forms a morphism of $H$-representations $W\to\res^G_HV$, by adjunction we get a $G$-morphism $\ind^H_GW\to V$.
This is an isomorphism if and only if $[G:H]\cdot\dim W=\dim V$.
And it is an injection iff the subspaces $(g_qW)_{q\in G/H}$ are linearly independent (i.e.\ have a direct sum in $V$).

\subsubsection*{Understanding $\Ind^H_G$}

Let us examine now $\Ind^H_G$.
Given a $k[H]$-module $M$, we can think of $\hom_{k[H]}(k[G],M)$ as a subspace of $\hom_k(k[G],M)$, which we can identify with $\St(G,M)$.
In particular, $\hom_{k[H]}(k[G],M)$ can be identified with maps $f\colon G\to M$ such that $f(hx)=hf(x)$ for all $h\in H$ and $x\in G$.
The $G$-action on this space is $(gf)(x)=f(xg)$ (as per the action on $\hom_{k[H]}(k[G],M)$).
So we can consider
$$ \Ind^H_G(M) = \set{f\colon G\to M}[\forall h\in H,x\in G.\,f(hx)=hf(x)],\qquad (gf)(x)=f(xg) $$

\subsubsection*{Back to business}

Note that $k[G]$ can be viewed as $\ind^1_G(k)$, and we can compute $[kG:E]$ using the adjunction $\ind^1_G\dashv\res^G_1$.

We want to compare $\ind^H_G$ and $\Ind^H_G$; we show that we have a natural transformation $\ind^H_G\to\Ind^H_G$, which is an isomorphism when $[G:H]$ is finite.
By the adjunction, a $G$-morphism $\ind^H_G(V)\to\Ind^H_G(V)$ is equivalent to an $H$-morphism $V\to\res^G_H\Ind^H_G(V)$ (we abuse notation and write $\Ind^H_G(V)$ for the right side).
We define such a morphism by mapping $v\in V$ to the function $G\to V$ defined by mapping $h\in H$ to $hv$ and $g\in G-H$ to $0$.

If $[G:H]$ is finite, we define its inverse $\Ind^H_G(V)\to\ind^H_G(V)$ as follows: let $(g_q)_{q\in G/H}$ be representatives of the cosets.
Then map $f\in\Ind^H_G(V)$ to $\sum_{q\in G/H}\delta_{g_q^{-1}}\otimes f(g_q)$.
This indeeds forms an inverse morphism.

So in our case, where $G$ is finite, we have that $\ind^H_G$ and $\Ind^H_G$ are isomorphic.

We can also abstract this argument.
Let $\iota\colon B\to A$ be a morphism of $k$-algebras.
Then we have the following natural morphism (in $M$):
$$ \hom_B(A,B)\otimes_B M \to \hom_B(A,M) $$
$\hom_B(A,B)$ is a right $B$-module since $B$ is a right $B$-module, so the left-hand side is well-defined.
And both sides are left $A$-modules since $A$ is.
This morphism maps $T\otimes m$ to $a\mapsto T(a)m$.

Now suppose $e\colon A\to B$ is a morphism of left $B$-modules.
Then we can construct a morphism of $(A,B)$-bimodules
$$ A\to\hom_B(A,B) $$
by mapping $a$ to $a'\mapsto e(a'a)$.
This gives a morphism of $A$-modules, natural in $M$:
$$ A\otimes_BM\to\hom_B(A,B)\otimes_BM $$

Composing gives us a morphism (natural in $M$):
$$ A\otimes_BM\to\hom_B(A,B)\otimes_BM\to\hom_B(A,M) $$
Note that this is a morphism $\ind^B_A\to\Ind^B_A$.

We claim that the second arrow is an isomorphism when $A$ is a free left $B$-module of finite rank.
Indeed, for a left $B$-module $N$, we can extend our second arrow to
$$ \hom_B(N,B)\otimes_BM\to\hom_B(N,M) $$
This is clearly an isomorphism when $N=B$.
And if it is an isomorphism for $N_1,N_2$ then it is an isomorphism for $N_1\oplus N_2$.
Thus if $N$ is a free left $B$-module of finite rank (so $N=B\oplus\cdots\oplus B$), then this is an isomorphism.

The first arrow is an isomorphism if $A\to\hom_B(A,B)$ is also an isomorphism.

In our case, we have $A=k[G]$, $B=k[H]\subseteq k[G]$ and $e\colon k[G]\to k[H]$ defined by $e(\delta_h)=\delta_h$ for $h\in H$ and $e(\delta_g)=0$ for $g\in G-H$.
Then $A$ is indeed a free left $B$-module of finite rank (when $[G:H]$ is finite): $k[G]=k[H]^{\oplus[G:H]}$.
And $A\mapsto\hom_B(A,B)$ is readily shown to be an isomorphism.

